Quatre càrregues elèctriques positives, d'$1,00\times 10^{-5}\ \mathrm{C}$ cadascuna, es troben als vèrtexs respectius d'un quadrat de $\sqrt{2}\ \mathrm{m}$ de costat. Calculeu:

\begin{apartat}{1}
L'energia necessària per a la formació del sistema de càrregues.
\end{apartat}

\begin{apartat}{1}
El valor de la càrrega elèctrica negativa que hem de situar al centre del quadrat perquè la força electrostàtica sobre cadascuna de les càrregues sigui nul·la.
\end{apartat}

\blocetiquetat{Dada:}{$k = 9,00\times 10^{9}\ \mathrm{N\,m^2\,C^{-2}}$}
